% =====================================================================
%  Thème 4 — Produit scalaire et orthogonalité — FACILE — Corrigé
% =====================================================================
\documentclass[11pt,a4paper]{article}
\usepackage[utf8]{inputenc}
\usepackage[T1]{fontenc}
\usepackage{lmodern}
\usepackage[french]{babel}
\usepackage{amsmath,amssymb,mathtools}
\usepackage{xcolor}
\usepackage{enumitem}
\usepackage{fancyhdr}
\usepackage[a4paper,margin=2.1cm,headheight=26pt,headsep=10pt]{geometry}

\definecolor{primary}{RGB}{222,70,80}
\definecolor{primarydark}{RGB}{150,30,42}
\definecolor{excol}{RGB}{0,135,90}

\pagestyle{fancy}\fancyhf{}
\fancyhead[L]{\small\textcolor{primarydark}{\textbf{Fiche 2} — Calcul vectoriel}}
\fancyhead[R]{\small\textcolor{primarydark}{\textsc{Thème 4 — Facile — Corrigé}}}
\fancyfoot[C]{\small\thepage}
\renewcommand{\headrulewidth}{0.8pt}

\newcommand{\vc}[1]{\overrightarrow{#1}}
\providecommand{\norm}[1]{\left\lVert #1\right\rVert}
\newcommand{\cor}[1]{\medskip\noindent{\color{primarydark}\bfseries\large Corrigé #1.}\par\smallskip}
\newcommand{\rep}{\textcolor{excol}{$\blacktriangleright$}\ }

\begin{document}
\begin{center}
{\LARGE\bfseries\color{primarydark}Produit scalaire et orthogonalité — Corrigé Facile}\\[2pt]
{\color{primary}\rule{0.55\linewidth}{1.1pt}}
\end{center}
\vspace{2mm}

\cor{1}
\begin{enumerate}[label=\textbf{\alph*)},leftmargin=2em,itemsep=2pt]
  \item $3\times1+2\times4=3+8=11$.
  \item $5\times2+(-1)\times3=10-3=7$.
  \item $(-2)\times3+4\times(-1)=-6-4=-10$.
  \item $2\times0+0\times5=0$.
\end{enumerate}
\rep On multiplie les premières composantes ensemble, puis les secondes, et on additionne.

\cor{2}
\begin{enumerate}[label=\textbf{\alph*)},leftmargin=2em,itemsep=2pt]
  \item $3\times(-2)+2\times3=-6+6=0$ : \textbf{orthogonaux}.
  \item $1\times2+4\times3=2+12=14\neq0$ : \textbf{non} orthogonaux.
  \item $6\times1+(-3)\times2=6-6=0$ : \textbf{orthogonaux}.
\end{enumerate}

\cor{3}
\begin{enumerate}[label=\textbf{\alph*)},leftmargin=2em,itemsep=2pt]
  \item $\vec u\cdot\vec u=3^2+4^2=9+16=25$, donc $\norm{\vec u}=\sqrt{25}=5$.
  \item $\vec u\cdot\vec u=1^2+(-2)^2=1+4=5$, donc $\norm{\vec u}=\sqrt5$.
\end{enumerate}
\rep $\vec u\cdot\vec u=\norm{\vec u}^2$ : la norme est la racine du carré scalaire.

\cor{4}
\begin{enumerate}[label=\textbf{\alph*)},leftmargin=2em,itemsep=2pt]
  \item $2\times3+1\times4=10>0$ : angle \textbf{aigu}.
  \item $2\times(-3)+1\times1=-5<0$ : angle \textbf{obtus}.
  \item $1\times2+2\times(-1)=0$ : angle \textbf{droit} (vecteurs perpendiculaires).
\end{enumerate}
\rep Le signe du produit scalaire suffit : $+$ aigu, $0$ droit, $-$ obtus.

\cor{5}
\begin{enumerate}[label=\textbf{\alph*)},leftmargin=2em,itemsep=2pt]
  \item $1\times2+2\times(-1)=0$ : les droites sont \textbf{perpendiculaires}.
  \item $3\times1+1\times3=6\neq0$ : les droites ne sont \textbf{pas} perpendiculaires.
\end{enumerate}
\rep Deux droites sont perpendiculaires ssi le produit scalaire de leurs vecteurs directeurs est nul.

\end{document}
